\chapterhead{3}{Falling and Floating}{Gravity, air and water push on everything}\label{ch3}

Let go of a book and it falls. Let go of a sheet of paper and it drifts down slowly, wobbling. Both are pulled down by gravity. So why is the paper so slow?

The answer is that something else acts on the paper. As a thing falls through air, the air pushes up against it. The more surface the thing spreads out, the bigger that upward push. Flat paper spreads a lot of surface, so the air pushes up on it almost as hard as gravity pulls down.

\begin{center}
\begin{tikzpicture}
  \fill[sky!30,draw=ink,line width=1pt] (-1.0,0) rectangle (1.0,0.12);
  \draw[-{Stealth},line width=2.4pt,color=tomato] (0,0.06) -- (0,-1.8); \node[font=\sffamily\scriptsize,anchor=west] at (0.15,-1.6) {gravity, 5};
  \draw[-{Stealth},line width=2.4pt,color=leaf!70!black] (0,0.06) -- (0,1.7); \node[font=\sffamily\scriptsize,anchor=west] at (0.15,1.5) {air, 4};
  \node[font=\sffamily\scriptsize] at (0,-2.4) {flat paper: leftover 1 down};
  \fill[tangerine!70,draw=ink] (5.0,0.4) circle (0.35);
  \draw[-{Stealth},line width=2.4pt,color=tomato] (5.0,0.4) -- (5.0,-1.4); \node[font=\sffamily\scriptsize,anchor=west] at (5.15,-1.2) {gravity, 5};
  \draw[-{Stealth},line width=2.4pt,color=leaf!70!black] (5.0,0.4) -- (5.0,0.9); \node[font=\sffamily\scriptsize,anchor=west] at (5.15,1.0) {air, 1};
  \node[font=\sffamily\scriptsize] at (5.0,-2.4) {crumpled paper: leftover 4 down};
\end{tikzpicture}
\end{center}

When the leftover is small, the paper speeds up only a little and falls slowly. When the leftover is big, the thing speeds up quickly. Squash the paper into a ball and the air pushes up much less.

\begin{tryit}{Flat or crumpled?}
\textbf{You need:} two sheets of the same scrap paper, a timer showing tenths of a second, a helper if you can.
\begin{enumerate}[leftmargin=1.6em]
\item Hold one flat sheet as high as you safely reach and let it go. Time the fall to the floor.
\item Crumple the other sheet into a tight ball. Drop it from the same height and time it.
\item Do five drops of each, in a mixed order. Find the average time and spread for each.
\end{enumerate}
\par\medskip
\renewcommand{\arraystretch}{1.6}
\begin{tabular}{|l|c|c|c|c|c|c|}
\hline
 & \textbf{1} & \textbf{2} & \textbf{3} & \textbf{4} & \textbf{5} & \textbf{Average (s)}\\ \hline
\textbf{Flat} & & & & & & \\ \hline
\textbf{Crumpled} & & & & & & \\ \hline
\end{tabular}
\end{tryit}
\adult{Do the drops on the floor, not from furniture, with nobody underneath.}

\begin{worked}{Paper drops (practice numbers)}
\textbf{Flat:} the ten drops from Book~1 averaged 2.0~s, with a spread of 0.6~s.\\
\textbf{Crumpled:} 0.7, 0.6, 0.7, 0.7, 0.6~s. Total $=3.3$. Average $3.3\div5=0.66$, about 0.7~s. Spread $0.1$~s.\\
The gap between averages is $2.0-0.7=1.3$~s, much bigger than either spread.
\end{worked}

\begin{multix}[hmm]\multixbot[hmm]\tcblower
Gravity pulls heavy things harder, so heavy things must always fall faster. I will drop a book and a ball of paper and write that the book wins. Is the book guaranteed to win?
\end{multix}

\subhead{Do heavy things fall faster?}
Drop a book and a crumpled ball of paper together from the same height, side by side. Listen. Do you hear one thump or two?

\begin{worked}{Book and paper ball (practice description)}
Dropped together from shoulder height, the book and the ball of paper landed so close together that we heard one thump. With a flat sheet instead of the ball, the book landed long before the paper. Spread-out paper catches the air, and the ball does not.
\end{worked}

The upward push of the air is small for the ball and the book, so gravity does almost all the work, and gravity pulls both down with the same speeding-up. The test cannot show gravity is exactly equal, only that any difference is too small to hear.

\trust{We dropped flat and crumpled paper from the same height, five times each, and compared the average fall times and the spread.}
{The gap, 1.3~s, was much bigger than either spread, so the shape of the paper made a real difference to the fall. Changing only the shape (same sheet of paper) points to the air as the cause.}
{That air is the only reason, or that other shapes follow the same pattern. A stopwatch read to a tenth of a second cannot show tiny differences. The one-thump test is not a measurement of gravity.}

\begin{keepbox}
Each arrow keeps the direction and size of a force. \textbf{We kept} \wline[3.4cm] \textbf{and we set aside} \wline[3.4cm].
\end{keepbox}

\clearpage

\subhead{Floating}
Water pushes up, too. A piece of clay in a ball sinks. Shape the same clay into a boat and it floats. The clay has not changed. What changed is the shape, which makes the water push up harder.

\begin{center}
\begin{tikzpicture}
  \fill[sky!25] (-1.6,-1.6) rectangle (4.6,0);
  \draw[line width=1.2pt] (-1.6,-1.6) rectangle (4.6,0.4);
  \fill[ink!60,draw=ink] (0,-1.2) circle (0.3);
  \draw[-{Stealth},line width=2.2pt,color=tomato] (0,-1.2) -- (0,-1.9); \draw[-{Stealth},line width=2.2pt,color=leaf!70!black] (0,-1.2) -- (0,-1.5);
  \node[font=\sffamily\scriptsize,anchor=north] at (0,-2.0) {ball: gravity bigger, sinks};
  \fill[ink!60,draw=ink] (3.0,-0.2) -- (2.4,0.1) -- (3.6,0.1) -- cycle;
  \draw[-{Stealth},line width=2.2pt,color=tomato] (3.0,-0.15) -- (3.0,-0.9); \draw[-{Stealth},line width=2.2pt,color=leaf!70!black] (3.0,-0.15) -- (3.0,0.6);
  \node[font=\sffamily\scriptsize,anchor=north] at (3.0,-2.0) {boat: water push equals gravity, floats};
\end{tikzpicture}
\end{center}

\begin{tryit}{Clay boats}
Make two clay boats from the same amount of clay: one wide and shallow, one narrow and deep. Float each on water. Add coins one at a time until it sinks. How many coins did each hold?\par\medskip
Wide boat: \wline[1.6cm] coins \qquad Narrow boat: \wline[1.6cm] coins
\end{tryit}
\adult{Clean up spilled water straight away, and keep coins and clay away from small children.}

\begin{worked}{Clay boat results (practice numbers)}
Same lump of clay for both. Wide, shallow boat: held 12 coins before sinking. Narrow, deep boat: held 6 coins. The wide boat pushed away more water, so the water pushed up on it harder.
\end{worked}

\begin{puzzle}{Time the drops}
A crumpled ball took 0.6, 0.7, 0.6, 0.8, 0.6~s to fall. A flat sheet took 1.9, 2.2, 2.0, 2.1, 1.8~s.\\
(a) Crumpled: total \wline[1.3cm] average \wline[1.3cm] \quad (b) Flat: total \wline[1.3cm] average \wline[1.3cm]\\
(c) Gap: \wline[1.3cm]~s \quad (d) Is the gap bigger than both spreads? \wline[1.6cm]
\end{puzzle}

\begin{puzzle}{Draw the arrows}
A leaf floats down. Draw gravity and the air push, with the air arrow a bit shorter than gravity.
\begin{center}
\begin{tikzpicture}
  \fill[leaf!60,draw=ink] (0,0) ellipse (0.6 and 0.15);
\end{tikzpicture}
\end{center}
Is the leaf speeding up or slowing down? \wline[3cm]
\end{puzzle}

\begin{puzzle}{Sink or float?}
For each, say what the pushes are and whether it sinks or floats: (a) a big wooden plank; (b) a coin; (c) a steel ship. Write a short reason for each.\\
\rule{\linewidth}{0.5pt}\\[4pt]\rule{\linewidth}{0.5pt}
\end{puzzle}

\begin{puzzle}{A fair test}
Marco tests whether paper shape changes fall time. He uses a big flat sheet and a small crumpled sheet from different drops. What else is different? \wline[4cm]\\
How would you fix it? \wline[4cm]
\end{puzzle}
\clearpage

\begin{puzzle}{How much longer?}
Flat paper took about 2.0~s to fall. Crumpled paper took about 0.7~s. How many seconds longer did the flat paper take? \wline[2cm]~s\\
About how many times as long? \wline[2cm]
\end{puzzle}

\begin{puzzle}{A parachute}
A toy parachute has a bigger canopy than a toy without one. The air pushes up on it with a bigger arrow. What will happen to the fall time? \wline[3cm]\\
Why? \wline[4cm]
\end{puzzle}

\begin{puzzle}{Wide and narrow boats}
The wide clay boat held 12 coins. The narrow boat held 6. How many times as many coins did the wide boat hold? \wline[2cm]\\
Why might a wider boat hold more? \wline[4cm]
\end{puzzle}

\begin{puzzle}{A speeding ball}
A ball is let go. At first gravity pulls down with size 5 and the air pushes up with size 1. What is the leftover? \wline[3cm]\\
Later, the ball falls at a steady speed. How big is the air's push now? \wline[2cm]
\end{puzzle}
\clearpage
